\section{Evaluation}
\label{sw-eval-section}

\paragraph{Materials and execution.}
We evaluated skill utility, governance, and source adaptation using Petstore/OpenAPI
creation, lookup, and deletion tasks. The materials comprised 60 task inputs,
60 governance states in 30 pairs, and 12 related source changes spanning routes,
required request fields, and enumeration values. Each change contained eight target
and eight retention inputs. Model-backed execution used deepseek-flash at
temperature~0, with zero automatic retries and output limits of 2,048 task tokens
and 4,096 proposal tokens. SkillWiki's local harness processed requests and supplied
responses without contacting the business API. The archived evaluation contains
4,590 unique model calls, including 18 update proposals; technical repetitions do
not increase the number of independent material units.

\paragraph{Skill utility.}
Each input was evaluated twice under three conditions. Table~\ref{sw-eval-utility}
shows that fixed skill representations achieved 120/120 correct answers, improving
on source material by 6.67 percentage points. Creation, lookup, and deletion each
achieved 40/40. Raw and final outputs coincided, attributing this observed benefit
to model use of the selected representations rather than output correction.

\begin{table}[t]
\centering
\caption{Skill utility on 60 inputs with two technical repeats. All answers completed and met the output format.}
\label{sw-eval-utility}
\begin{tabular}{lrr}
\hline
Condition & Correct & Accuracy \\
\hline
No domain material & 84/120 & 70.00\% \\
Source material & 112/120 & 93.33\% \\
Fixed skill & 120/120 & 100.00\% \\
\hline
\end{tabular}
\end{table}

\paragraph{Evidence-based governance.}
The governance states exercised valid release, identity/content binding,
applicability, completeness, deferral/resolution, and recovery/state consistency.
An independent Git/CI+ implementation applied the same evidence contract and
performed its own repository actions. Both routes made correct decisions on 50/50
main states and 10/10 remaining states at the supplementary revision, with no false
accepts or state inconsistencies. One SkillWiki input-validation error was within
its preregistered allowed decision set and retains that classification. Agreement
with this explicitly implemented Git/CI+ workflow demonstrates the integration of
its evaluated governance requirements into SkillWiki's native skill lifecycle.

\paragraph{Source adaptation and lifecycle.}
The 12 events compared an original-rule reference, an unchanged skill, and an
updated skill, with two repeats: 1,152 core answers. Under changed requirements,
target correctness increased from 0/192 for unchanged skills to 176/192 raw model
answers and 192/192 final outputs; explicit source checks corrected 16 answers.
Retention was 173/192 against the unchanged baseline's 192/192, with seven events
meeting the joint improvement-and-retention criterion. Full qualification used
36 cases per identity. Five updated events per route qualified, were released,
and completed recovery with revalidation and readback in the main evaluation.

In subsequent revisions separating request and response constraints, five revised
candidates achieved 180/180 final qualification answers (171/180 raw). Their five
recovered identities achieved 180/180 final answers (167/180 raw), rejected old
proofs, and were revalidated and republished. These candidates comprise four R1
identities and one R2 identity; U06 R1 remains recorded at 35/36. Supplementary U08
updated and recovered identities each achieved 36/36 raw and final answers and
completed release/readback. Across recorded revisions, 11 events have successful
lifecycle chains. U03 remains unpublished at 35/36 because an out-of-range integer
is rejected with an error category that differs from the contract.

\paragraph{Scope and artifacts.}
The results establish utility for selected skills, agreement with a carefully
implemented governance baseline, and executable evidence-gated adaptation and
recovery. The materials comprise development/regression cases and related
extensions within one API family. Results retain their runtime and candidate
revisions; cross-revision coverage is not a single-version success rate. The
public companion provides synthetic inputs, expected answers, source snapshots,
governance states, result tables, and a de-identified call index.
